\chapter{Future work}

The advanced fluid neutral boundary conditions should be kept up to date with any advancements in plasma boundary modeling, and further extensions for application to more realistic cases should be considered. Examples of possible extensions are:

\begin{itemize}
    \item \textbf{Multi-species cases:} mixtures of H, D, or T have not yet been considered. Both for the transport model and the boundary conditions, the theory needs to be revisited for multi-species cases. For the boundary conditions, we suspect that little or no difficulty arises when using the truncated Maxwellian approximation for the incident atoms. However, when considering the diffusion approximation, the collisions between the different species must be taken into account.
    \item \textbf{Impurity species:} the same principles that were used here for the hydrogenic species could possibly lead to improved boundary conditions for the fluid modeling of impurity neutrals. It would then again be advised to use the truncated drifting Maxwellian assumption instead of the diffusion approximation for the incident atoms. For the recycling, separate contributions for each ionization stage must be considered, taking also into account the differences in sheath acceleration depending on the ionization stage.
    \item \textbf{Advanced sheath models:} we have used a simple sheath acceleration model, corresponding to what EIRENE does. If more advanced sheath acceleration models become available in EIRENE, the recycling part of the AFN boundary conditions can be improved accordingly.
    \item \textbf{The role of perpendicular bulk velocities:} in the end, contributions from smaller perpendicular bulk velocities (both of the incident ions and incident atoms) were neglected to reduce the dimensionality of the pre-integrated databases. The effect of this simplification is not yet fully understood. However, good fluid-kinetic agreement was achieved in several AFN publications, including one where drifts were included in the plasma model \cite{psi_wim}. Still, a more complete treatment of all velocity components could be desired in the future.
    \item \textbf{Neutral-neutral collisions:} when extending the fluid neutral models to include neutral-neutral collisions, the AFN boundary conditions must again be revisited. When using the truncated drifting Maxwellian assumption for the incident atoms, likely little or no change will be required. When using the diffusion approximation, the atom-atom collisions must be included in the derivation.
    \item \textbf{New reflection models:} if new reflection databases are released as a replacement for TRIM, the AFN boundary conditions should be extended to include those. Depending on the precise format of those new databases, most elements of the current procedure can be reused.
    \item \textbf{Fluid models for molecules:} if an approximate fluid description for the molecules is developed, the AFN boundary conditions can be easily adapted to provide realistic boundary conditions for these fluid molecules. The fraction of incident ions and atoms that is thermally released as molecules is already calculated. What needs to be added is a model for the incident molecules that return as thermally released molecules.
\end{itemize}

\noindent When changing or extending the AFN boundary conditions, please also update this document. The \LaTeX~source files can be found in {\tt \$SOLPSTOP/doc/AFN\_documentation}.
